\section[Notebook 1972-2: splitting finite-rank endomorphisms; an integral over a triangle]{Notebook 1972-2: splitting an endomorphism of a free module of infinite rank; an integral over a triangle\protect\verifie{Notebook1972_2}}\label{sec:quillen-1972-2}

\emph{Context.} In February 1972, \carnet{1972}{2} looks for a category with the homotopy type of $BGL(k)^+$ and studies what happens on passing to an infinite-dimensional space $W$. For a free module $R^{(S)}$ on an infinite set $S$, Quillen considers the group $\operatorname{Aut}^f(R^{(S)})$ of automorphisms $1 + K$ with $K$ of finite rank, and conjectures (page~18) that the inclusion $GL_\infty(R) \subset \operatorname{Aut}^f(R^{(\infty)})$ induces an isomorphism on cohomology. Over a field, every such $1 + K$ comes from an automorphism of a finite-dimensional summand (page~17). The entry of 19 February asks how far this holds over a ring. The second statement belongs to the last block of the notebook, on the Godbillon--Vey class and a $2$-cocycle on the group of compactly supported diffeomorphisms of $\mathbf R$; on page~46, Quillen computes an elementary integral on the way to the cocycle of pages~49--50.

\emph{The page (pp.~20--21).} Let $R$ be a ring and $K$ an endomorphism of $R^{(S)}$ of finite rank, that is, with image in $R^T$ for a finite $T \subset S$. Quillen asks when there is a decomposition $R^{(S)} = V \oplus Q$ with $V$ finitely generated, $K(V) \subset V$ and $K(Q) = 0$; then $1 + K$ is the automorphism $(1 + K|_V) \oplus 1$, which comes from $\operatorname{Aut}(V)$. If $\operatorname{Im} K$ is finitely generated, he picks basis vectors $e_i$, $i \in T_1$, whose images generate it, writes $K e_j = \sum_{i \in T_1} r_{ji} K e_i$ for the other $j$, and changes coordinates by the vectors $e_j - \sum_i r_{ji} e_i$, which lie in the kernel; the span of $e_j$, $j \in T_1 \cup T$, and the span of the new vectors with $j \notin T_1 \cup T$, then do the job. Conversely, $\operatorname{Im} K = K(V)$ is finitely generated. He concludes that over a non-noetherian ring some finite-rank maps $R^{(S)} \to R$ have kernels containing no summand of finite corank, and on page~22 passes to a noetherian ring.

\emph{Intuition.} A summand $Q$ in the kernel of $K$ with a finitely generated complement $V$ forces $\operatorname{Im} K = K(V)$ to be finitely generated. Conversely, once a finite part $R^{(U)}$ of the basis already has the full image under $K$, every vector outside $R^{(U)}$ can be corrected by a vector of $R^{(U)}$ with the same image, and the corrected vectors span a complement in the kernel. The finite-rank hypothesis is not needed: a finitely generated image lies in some $R^T$ automatically.

\begin{encadre}
\begin{theoreme}[after Quillen, notebook 1972-2, pp.~20--21]\label{thm:quillen-1972-2}
Let $R$ be a ring, $S$ a set and $K$ an endomorphism of the free left module $R^{(S)}$. There is a decomposition $R^{(S)} = V \oplus Q$ with $V$ finitely generated, $K(Q) = 0$ and $K(V) \subset V$ if and only if $\operatorname{Im} K$ is finitely generated; $V$ can then be taken of the form $R^{(U)}$ with $U \subset S$ finite. In particular, if $R$ is left noetherian, every endomorphism of finite rank splits in this way; if $R$ is not left noetherian and $I \subset R$ is a left ideal that is not finitely generated, the map $R^{(I)} \to R$, $e_i \mapsto i$, has no such $Q$ in its kernel.
\end{theoreme}
\end{encadre}

\noindent The statement drops the page's hypothesis that $K$ is of finite rank, and makes the page's remark on non-noetherian rings explicit with the ideal $I$.

\begin{proof}
If $R^{(S)} = V \oplus Q$ with $V$ finitely generated and $K(Q) = 0$, then $\operatorname{Im} K = K(V) + K(Q) = K(V)$ is finitely generated. Conversely, let $g_1, \dots, g_n$ generate $\operatorname{Im} K$ and choose $x_k$ with $K x_k = g_k$. Let $U \subset S$ be the finite union of the supports of the $g_k$ and $x_k$, $V = R^{(U)}$ and $W = R^{(S \setminus U)}$, so that $R^{(S)} = V \oplus W$, $\operatorname{Im} K \subset V$ (hence $K(V) \subset V$) and $\operatorname{Im} K = K(V)$. The map $K|_V \colon V \to \operatorname{Im} K$ is onto and $W$ is free, so $K|_W \colon W \to \operatorname{Im} K$ lifts to $h \colon W \to V$ with $K h = K|_W$. Put $Q = \{w - h(w) : w \in W\}$. Then $K(Q) = 0$; if $w - h(w) \in V$ then $w \in V \cap W = 0$; and $v + w = (v + h(w)) + (w - h(w))$, so $R^{(S)} = V \oplus Q$. If $R$ is noetherian and $\operatorname{Im} K \subset R^T$ with $T$ finite, $\operatorname{Im} K$ is finitely generated. For the last assertion, the map $e_i \mapsto i$ has image $I$, which is not finitely generated.
\end{proof}

\noindent The lift $h$ replaces the page's explicit change of coordinates: with $h(e_j) = \sum_i r_{ji} e_i$ one recovers the vectors $e_j - \sum_i r_{ji} e_i$.

\emph{The page (p.~46).} In the course of the cocycle computation, Quillen integrates $(1 + ax + by)^{-2}$ over the triangle $x, y \ge 0$, $x + y \le 1$, first in $y$, then in $x$, and finds the symmetric expression below; the function $\log(1+x)/x$ appears in the cocycle of page~50.

\begin{encadre}
\begin{theoreme}[Quillen, notebook 1972-2, p.~46]\label{thm:quillen-1972-2-integral}
For real numbers $a, b > -1$ with $a \ne 0$, $b \ne 0$ and $a \ne b$,
\[
\int_0^1 \!\! \int_0^{1-x} \frac{dy\,dx}{(1 + ax + by)^2} = \frac{\log(1+a)}{a(b-a)} + \frac{\log(1+b)}{b(a-b)}.
\]
\end{theoreme}
\end{encadre}

\noindent The page states no hypothesis; $a, b > -1$ is the condition for $1 + ax + by$ to be positive on the closed triangle, whose vertices give the values $1$, $1 + a$, $1 + b$, and the other three exclude the zero denominators of the formula.

\begin{proof}
On the triangle, $1 + ax + by = (1 - x - y) + x(1 + a) + y(1 + b) > 0$. For $0 \le x \le 1$, an antiderivative in $y$ is $-1/(b(1 + ax + by))$, so the inner integral is $\frac1b\bigl(\frac{1}{1+ax} - \frac{1}{1 + b + (a-b)x}\bigr)$. An antiderivative of this in $x$ is $\frac1b\bigl(\frac{\log(1+ax)}{a} - \frac{\log(1 + b + (a-b)x)}{a-b}\bigr)$, and its variation from $0$ to $1$ is $\frac{\log(1+a)}{ab} - \frac{\log(1+a) - \log(1+b)}{b(a-b)}$, which is the right-hand side since $\frac{1}{ab} - \frac{1}{b(a-b)} = \frac{1}{a(b-a)}$.
\end{proof}

\noindent The proof is the computation of the page, with the positivity made explicit.

\emph{Verification in Lean.} The file \href{\leanblobmain/Notebook1972_2.lean}{\texttt{Notebook1972\_2.\allowbreak{}lean}} writes $R^{(S)}$ as \texttt{S →₀ R} and $R^{(U)}$ as \texttt{Finsupp.supported R R U}, over any \texttt{Ring R}. The lifting step, for any decomposition $M = V \oplus W$ with $W$ projective and $K(V) = K(M)$, is \texttt{exists\_isCompl\_le\_ker}; the converse is \texttt{fg\_range\_of\_isCompl}. Theorem~\ref{thm:quillen-1972-2} is \texttt{exists\_splitting\_iff}, with the noetherian case \texttt{exists\_splitting\_of\_isNoetherianRing} and the counterexample \texttt{not\_exists\_splitting\_of\_not\_isNoetherianRing}. Theorem~\ref{thm:quillen-1972-2-integral} is \texttt{integral\_simplex}, stated with iterated interval integrals, the inner one being \texttt{integral\_inner}; the integral over the triangle as a subset of $\mathbf R^2$ for Lebesgue measure, which equals the iterated one by Fubini's theorem, is not formalized. The axioms used are the three usual axioms of mathlib.

\litterature{the splitting lemma was not found in this form; it is elementary, and was not searched for specifically. The integral is elementary calculus; the formula was checked numerically for $a = 0.7$, $b = 2.3$. The edition presents neither as new.}
